\section{The whole bare component}\label{ms:components}
An estimate of a tensor with one high incidence cannot generally use
the complete tensor at that same index. The reflection in
Lemma~\ref{ms:reflected} changes the pairing across that edge. We
now keep that change throughout the neighboring component.

\begin{theorem}[Component comparison]\label{ms:component}
Let $A$ be a fixed unturned bare component, and $Q=\Turn(A)$ its
typed output. Suppose its incoming full coefficient arrays obey
\begin{equation}\label{ms:incoming}
 |f_{j,r}|\le e^{N(c_j+o(1))}
 \begin{cases}e^{Ng(r/d)},&r\le m,\\1,&r>m,\end{cases}
\end{equation}
uniformly on a stated meridian compact. Let $A^0$ be its original
complete positive reference with ordinary boundary profiles $g$,
and put $\Omega=\Phi_{C(hA^0)}+\sum_jc_j$, using the actual output
probe of \cite[Theorem 7.4]{Newton}. Then
\[
 Q=Q^{\mathrm{clean}}+E,\qquad
 \normJ Q\le e^{N(\Omega+o(1))},\qquad
 \normJ E\le e^{N(\Omega-\eta+o(1))}
\]
for some $\eta>0$. In the clean term all bare tensors are complete,
all original external inputs and the output are low, and every
complete-supported internal index is retained. Actual low arrays,
physical edge weights and typed scalars are retained. A loss in one
entire incoming array is inherited without change.
\end{theorem}
\begin{proof}
Expand each moving vertex as $V_\sigma=V_0+\delta V_\sigma$, then
partition every index into its low and high bands. There are only
finitely many patterns. Unselected and constant vertices retain their
positive complete tensors. The complete pair support permits at most
one high incidence at each such vertex.

For a selected difference use the all-low estimate
of Lemma~\ref{ms:low-trace}, the one-high reflected estimate of
Lemma~\ref{ms:reflected}, or the full norm bound
\begin{equation}\label{ms:several-high}
 |\delta\phi_\sigma(\prod_i e_{n_i})|
       \le N^A\prod_{n_i\le m}w_{n_i}.
\end{equation}
The last inequality follows from \eqref{ms:basis-norm}, fixed-valence
quotient reduction and the polynomial trace norm. It retains the
high normals, even when all their values vanish.

Consider a selected one-high difference whose high incidence is
internal. Its neighboring vertex is constant because moving
equal-color vertices are nonadjacent. Each endpoint has only this
one high incidence. Hence all such selected edges form a matching
$S$. Reflect $h$ to $r=d-h$ at the moving endpoint and use
\eqref{ms:fusion}. Summing the complementary pair fuses the two
whole vertices in their induced order. Removing its band restriction
only enlarges the positive sum. Every fused vertex has valence at
least four and low exterior incidences. The restored complementary
label remains one hidden unturned wire in a positive resolution.
No other branch is replaced by an independent maximum. By
Lemma~\ref{ms:deletion} the resulting pressure is at most the original
pressure. For a one-high external incidence there is no fusion:
reflection makes its array rate $c_j$ while the low reference has
rate $c_j+g(r/d)$. Keep the nonnegative surplus $g(r/d)$.

Selected vertices with several high incidences are not fused.
They remain nonadjacent. Remove them using \eqref{ms:several-high},
assigning one $g$ to each low arm and marking every high arm in the
retained old region. Mark also an unreflected high original input
or output. Every retained whole vertex has at most one mark; the
matching fusions did not create high exterior incidences.
The reassembly following Lemma~\ref{ms:half-boundaries} moves these
marks simultaneously to half and reinserts the removed vertices.
Each charged $g$ is consumed once by their cardinal action. All
original constants remain once. The result is a boundary state of
the same fused complete tree: either a pair sum is one or an
ordinary external leaf is half. Its unique strict maximum avoids
this closed set, so it has a fixed loss. This includes a component
consisting of a single removed vertex. With no removed vertex the
same half-boundary argument gives the high-input or high-output
loss directly.

It remains to consider a nonempty difference pattern with no such
boundary. Its positive comparison is a fused tree $\mathcal G/S$.
Outside a small neighborhood of its unique strict resolved maximum,
compactness gives a fixed loss. Inside that neighborhood there are
three possibilities. An externally reflected label lies strictly
between zero and half, so its saved $g$ is strictly positive. If
an internal reflection occurs, the fused tree has a vertex of
valence at least four. It cannot be the exceptional entire
two-trivalent-vertex tree of Lemma~\ref{ms:deletion}. Every physical
exterior label of that fused pair is therefore in a compact low
interval. Its restored wire $r/d$ is strictly positive. A maximum
above half is separated from the sector $r\le m$; otherwise the
strict part of Lemma~\ref{ms:reflected} applies, including at a
limiting half wire. Finally, with no reflection, every selected
difference is all-low in the original tree. The strict whole-vertex
support gives a positive total slack. Lemma~\ref{ms:regular-trace}
then gives a loss. In the two-trivalent-vertex exception the raw
cubic degree on the entire strict neighborhood is less than
$2N-1$, so its trace difference is exactly zero instead.

These cases exhaust the nonempty difference patterns. One strict
factor is enough; all other factors retain their polynomial
majorants. Uniform factorial estimates on the closed support
polytope give $O(\log N)$ errors and only polynomially many states.
Take the minimum of the finitely many positive losses and decrease
it once to absorb those errors.

For the output use the exact identity \eqref{ms:turned-probe}:
\[
 \normJ Q\le C N^A\sum_{r=0}^d
                     w_r|\phi_\sigma(e_r A)|.
\]
This is precisely the ordinary low probe $g$ and the zero-rate high
probe already treated. Its action is $\Omega$, by
\cite[Theorems 7.3--7.4]{Newton}; it is not a newly inserted
physical circle. In the empty difference pattern with low external
indices, the positive complete coefficients and bounded edge
weights give the original positive reference, including its full
low-output profile. Every internal complete-supported index stays
in that term. This proves the asserted decomposition. Multilinearity
preserves a loss in one full array in all comparisons, and shifts
the corresponding constant $c_j$ without changing any optimizer.
\end{proof}

\begin{figure}[t]
\centering
\begin{tikzpicture}[scale=.85,>=Latex,every node/.style={font=\small}]
\node[circle,draw,minimum size=10mm] (m) at (0,0) {$\delta V_\sigma$};
\node[circle,draw,minimum size=10mm] (c) at (2.5,0) {$V_0$};
\draw[thick] (m)--node[above]{$h,h$}(c);
\draw (m)--(-.9,.9) node[above]{$a_1$};
\draw (m)--(-.9,-.9) node[below]{$a_2$};
\draw (c)--(3.4,1) node[above]{$b_1$};
\draw (c)--(3.6,0) node[right]{$b_2$};
\draw (c)--(3.4,-1) node[below]{$b_3$};
\draw[->,thick] (4.5,0)--node[above]{$r=d-h$}(6,0);
\node[draw,ellipse,minimum width=22mm,minimum height=15mm] (f) at (7.5,0) {$V_0^{\mathrm{fused}}$};
\draw (f)--(6.4,1) node[above]{$a_1$};
\draw (f)--(6.4,-1) node[below]{$a_2$};
\draw (f)--(8.7,1) node[above]{$b_1$};
\draw (f)--(9.1,0) node[right]{$b_2$};
\draw (f)--(8.7,-1) node[below]{$b_3$};
\node at (7.5,-1.8) {$\displaystyle\sum_{r=0}^{d}V_0(\mathbf a,r)V_0(d-r,\mathbf b)$};
\end{tikzpicture}

\caption{A selected high incidence joins a moving trace (left disk)
to a constant trace. The original covariant labels agree, $h,h$.
Reflection changes only the left label to $r=d-h$. The complete
metric then pairs $r$ with $d-r$ and fuses the disks in their
boundary order. Every other arm remains in the same contraction.
The picture depicts the tensor calculation used for an upper bound.}
\label{ms:reflection-figure}
\end{figure}

We explain how this component statement is used in the original
tree. A circle vertex, including its fixed signed crossing and
typed turn, converts the actual incoming spectral norm into
\eqref{ms:incoming} by Lemma~\ref{ms:coupon-bound}. A maximal bare
component uses Theorem~\ref{ms:component} and preserves its full
clean low profile. Thus a small spectral-norm error is converted
to a coefficient cone at a circle interface, where the exact
extractor proves that conversion. It is never converted by an
unsupported inversion of the Newton coefficient matrix inside a
bare component. The mutual induction terminates because the
original reduced graph is a finite tree. The root budgets of
\cite[Corollary 7.5]{Newton} identify its accumulated pressure
with $F_0$, retaining every additive constant.

This also gives fresh angular localization at any exposed circle.
The clean output-probe profile is $f_A(1-u)+g(u)+\sum c_j$, with a
unique maximizing $u\in(0,1/2)$. Outside a compact interval about
it there is a strict loss. Within it the low Newton angular estimate
\cite[Lemma 7.2]{Bounds} gives a negative quadratic angular term.
Errors have a strict spectral-norm gap, so complete Hermite
conditioning and the degree-$d$ Bernstein bound absorb them on a
sufficiently small complex strip. This uses the actual accumulated
pressure and works for an entirely bare-rooted expression by
exposing any original circle through the typed adjunction. For
$H=2$ the same exact pairing is the two-leaf pure source; no
nonexistent bare component is inserted.
